\begin{abstract}
This thesis investigates hybrid classical--quantum transfer learning
for galaxy morphology classification. A balanced dataset of 5,628
SDSS images, labelled using Galaxy Zoo classifications, supports a
binary task distinguishing elliptical and spiral galaxies and a
three-class task that additionally includes systems labelled as mergers.
A frozen, ImageNet-pretrained ResNet18 provides a shared representation
for classical heads and hybrid heads with four or eight simulated qubits.
The main comparison pairs the trainable quantum head with a classical
bottleneck of the same dimension; a frozen-circuit variant assesses the
contribution of training the circuit parameters.

The evaluation comprises 640 runs across ten stratified holdout
repetitions and four training regimes: 100, 300, 900 images and the
complete training split. Nested training subsets and shared validation
and test sets support paired comparisons. Macro-F1, bootstrap confidence
intervals and Holm-adjusted tests are complemented by per-class analysis,
probabilistic metrics and computational costs.

On the full binary task, the models achieve a Macro-F1 close to 0.96,
with no statistically supported advantage for the quantum head over the
matched classical control. On the full three-class task, the quantum
head has a lower mean Macro-F1 by 0.73 and 1.67 percentage points with
four and eight qubits, respectively. Reduced-data regimes do not reveal
a quantum advantage under the adopted protocol, while simulation requires
substantially more training time. The study provides a reproducible
benchmark that distinguishes predictive performance from the added
value of the quantum component and identifies concrete directions for
further investigation. Its conclusions concern the tested architectures
and simulation conditions, not quantum machine learning in general.
\end{abstract}
